\begin{abstract}
Small EEG models trained from scratch keep beating foundation models with 20 to 40 times more parameters on clinical event and seizure corpora, and the field has no account of why. We give one, and we test it. Clinical labels are defined by neurophysiologists per frequency band in terms of two quantities: how a band's amplitude envelope evolves (periodic discharges, artifact, seizure evolution) and what shape its waveform has (spike-and-wave). We formalize this as the \emph{demodulation hypothesis}: a clinical label is a low-complexity functional of the per-band log-envelope and band-limited waveform at the event timescale; a tokenizer that computes this grid explicitly raises the information a fixed-capacity encoder can use (predictive $\mathcal{V}$-information), and an encoder whose attention factorizes along time, electrode and band reads it with far fewer degrees of freedom than flat attention. The hypothesis yields eight falsifiable predictions, which we test with controlled signals and with 425 single-field contrasts drawn from 748 archived training runs on twelve clinical corpora. The large, replicated effects are the predicted ones: one learned band instead of eight costs 0.42 and 0.11 AUC-PR on CHB-MIT and TUSZ, a flat encoder costs 0.53 and 0.22, the envelope view alone matches the full model on TUEV, removing the waveform view costs 0.17 on CHB-MIT, and sixteen learned rows of the raw patch trail sixteen physical bands by 0.135 $\kappa$ under the same encoder. Content transforms that carry no label information behave as predicted too: an isometric rotation of the token by a measured cross-frequency direction changes the score by 0.003 on average over eight paired runs in the factorized reader, moves a flattened host reader by up to 0.2 in either direction, and costs up to 0.24 when the direction is content-dependent. The resulting 1.6M-parameter model reaches TUEV $\kappa$ 0.693, TUSZ AUC-PR 0.671 to 0.680 and IIIC $\kappa$ 0.497 against 0.575, 0.545 and 0.436 for the strongest reproduced baselines, and we report CHB-MIT on its natural window grid, on which every model scores lower than under the enriched evaluation used before.
\end{abstract}
